\documentclass[11pt]{article}

\usepackage[margin=1in]{geometry}
\usepackage{amsmath,amssymb,amsthm,mathtools}
\usepackage{booktabs}
\usepackage[hidelinks]{hyperref}

\newtheorem{theorem}{Theorem}[section]
\newtheorem{lemma}[theorem]{Lemma}
\newtheorem{proposition}[theorem]{Proposition}
\newtheorem{corollary}[theorem]{Corollary}
\newtheorem{definition}[theorem]{Definition}
\newtheorem{remark}[theorem]{Remark}

\DeclareMathOperator{\End}{End}
\DeclareMathOperator{\Aut}{Aut}

\hypersetup{
  pdftitle={Cores of Binary 3-Qualitative Independence Hypergraphs},
  pdfauthor={Jihun Kim},
  pdfsubject={Hypergraph cores and qualitative independence},
  pdfkeywords={qualitative independence, hypergraph core, covering array,
    hypergraph homomorphism, merged Johnson graph, Steiner system}
}

\title{Cores of Binary \(3\)-Qualitative Independence Hypergraphs}
\author{Jihun Kim\\
\small \texttt{jihunkimkw@gmail.com}}
\date{30 August 2026}

\begin{document}
\maketitle

\begin{abstract}
Let \(3\text{-}\mathrm{QI}(n,2)\) be the \(3\)-uniform hypergraph whose
vertices are unordered bipartitions of an \(n\)-element set into two blocks
of size at least \(4\).  Three such bipartitions form an edge if their common
refinement has all eight cells nonempty.  We prove that this hypergraph is a
core for every \(n\ge8\), resolving an open problem of Thomas and Akhtar.
Our main structural result is a lift theorem: if the most balanced
layer is a core, then so is the full hypergraph.  Its proof shows that
endomorphisms preserve the balanced layer and uses a separation argument to
extend rigidity from that layer to every vertex.  We then prove that the
balanced layer is a core for every even \(n\).  For
\(n\equiv2\pmod4\), a symmetry reduction and a divisibility
obstruction to a Steiner system exclude the only possible proper core.  For
\(n\equiv0\pmod4\) and \(n\ge16\), link-clique bounds together with explicit
combinatorial constructions rule out noninjective endomorphisms; \(n=12\) is
handled by a separate certificate, while \(n=8\) is the known base case.
Combining the even result with the odd balanced-layer
theorem proved in a companion manuscript gives the result for all \(n\ge8\).
\end{abstract}

\medskip
\noindent\textbf{Keywords.}
Qualitative independence; hypergraph core; hypergraph homomorphism; covering
array; merged Johnson graph; Steiner system.

\medskip
\noindent\textbf{MSC 2020.}
05C65, 05C60, 05B15, 05D05, 68R05.

\section{Introduction}

A finite hypergraph is a \emph{core} if every endomorphism is an
automorphism; see \cite{HellNesetril2004} for the graph setting.  In a binary
qualitative-independence hypergraph, a vertex is an unordered bipartition of
an \(n\)-element ground set, and three vertices form an edge when their common
refinement has all eight cells nonempty.  These hypergraphs encode triples of
columns in binary covering arrays
\cite{MeagherStevens2005,AkhtarMaity2017}.

Write \(H_n=3\text{-}\mathrm{QI}(n,2)\), and let \(B_n\) be the subhypergraph
induced by the most balanced bipartitions.  Thus \(B_n\) is
\(3\text{-}\mathrm{UQI}(n,2)\) when \(n\) is even and
\(3\text{-}\mathrm{AUQI}(n,2)\) when \(n\) is odd.  Thomas and Akhtar proved
that \(H_8=B_8\) is a core and asked for the core of \(H_n\) for \(n>8\)
\cite{ThomasAkhtar2026}.  They also identified the \(2\)-sections of the
balanced layers with merged Johnson graphs.

For odd \(n\ge9\), the companion work \cite{KimMergedJohnson2026} determines
the endomorphism monoid of the corresponding merged Johnson graph and thereby
shows that the balanced layer \(B_n\) is a core.  The present paper treats the
even balanced layers and establishes the passage from \(B_n\) to \(H_n\).

Our lift theorem shows that endomorphisms preserve the balanced layer and that
rigidity on this layer extends to the full hypergraph.  Maximum link cliques
identify the balanced layer intrinsically, while a separation lemma forces an
endomorphism that fixes the layer pointwise to fix every vertex.  For the even
balanced layers, we analyse their \(2\)-sections, using a Steiner-system
obstruction in one congruence class and a direct collision argument in the
other.  Together
with the odd balanced-layer theorem, this resolves the problem of Thomas and
Akhtar for every \(n>8\).

\begin{theorem}\label{thm:main}
For every integer \(n\ge8\), the hypergraph \(H_n\) is a core.
\end{theorem}

\begin{theorem}[Even balanced layers]\label{thm:even-middle}
For every even integer \(n\ge8\), the balanced layer
\(B_n=3\text{-}\mathrm{UQI}(n,2)\) is a core.
\end{theorem}

\begin{theorem}[Balanced-layer lift]\label{thm:lift}
Let \(n\ge8\).  If the balanced layer \(B_n\) of \(H_n\) is a core, then
\(H_n\) is a core.
\end{theorem}

\section{Qualitative independence and cores}

Let \(\Omega\) be an \(n\)-element set.  A binary string is identified, up to
interchanging its two symbols, with an unordered bipartition
\[
[A]=\{A,\Omega\setminus A\}.
\]
For a subset \(A\) of the current ground set, write \(A^c\) for its
complement.  Thus, for the present set \(\Omega\),
\(A^c=\Omega\setminus A\).  Write \(A^1=A\) and \(A^0=A^c\).  When an
orientation has been fixed, we occasionally denote the vertex \([A]\) by its
chosen representative \(A\).  Throughout, a homomorphism
of simple \(3\)-uniform hypergraphs is a vertex map taking every \(3\)-edge
to a \(3\)-edge; an endomorphism has the same source and target.
In the standard definition of \(3\text{-}\mathrm{QI}(n,2)\), each binary
symbol occurs at least \(4=2^{3-1}\) times; equivalently, both blocks have
size at least \(4\) \cite{ThomasAkhtar2026}.  Three vertices
\([A],[B],[C]\) form an edge when all eight cells
\[
A^{\epsilon_1}\cap B^{\epsilon_2}\cap C^{\epsilon_3},
\qquad
(\epsilon_1,\epsilon_2,\epsilon_3)\in\{0,1\}^3,
\]
are nonempty.  This condition is independent of the chosen orientations.

For \(n\ge8\), the \emph{balanced layer} (also called the \emph{middle
layer}) \(B_n\) is the subhypergraph induced
by the bipartitions with block
sizes
\[
\left\lfloor\frac n2\right\rfloor,
\qquad
\left\lceil\frac n2\right\rceil.
\]
It is \(3\text{-}\mathrm{UQI}(n,2)\) for even \(n\) and
\(3\text{-}\mathrm{AUQI}(n,2)\) for odd \(n\).

The \(2\)-section \([H]_2\) of a \(3\)-uniform hypergraph \(H\) is the
graph on \(V(H)\) in which two vertices are adjacent when they lie together
in a hyperedge.

\begin{lemma}[Graph-to-hypergraph transfer]\label{lem:transfer}
Let \(H\) be a finite simple \(3\)-uniform hypergraph.  If \([H]_2\) is a
core, then \(H\) is a core.
\end{lemma}

\begin{proof}
Every endomorphism of \(H\) is an endomorphism of \([H]_2\).  If
\([H]_2\) is a core, that map is injective.  On a finite vertex set, it is
bijective.  The induced map on \(3\)-subsets is then a permutation, and it
maps the finite edge set injectively into itself, hence onto itself.
\end{proof}

\section{Crossing cuts and the balanced-layer lift}

Throughout this section, \(n\ge8\).

For a set \(T\) of size \(t\), let \(\mathrm{Cross}(t)\) be the graph whose
vertices are nontrivial unordered bipartitions of \(T\), two of which are
adjacent when all four cells of their common refinement are nonempty.

\begin{proposition}\label{prop:cross}
For \(t\ge4\),
\[
\omega(\mathrm{Cross}(t))
=q(t):=\binom{t-1}{\lfloor t/2\rfloor-1}.
\]
Moreover, \(q(t+1)>q(t)\) for every \(t\ge4\).
\end{proposition}

\begin{proof}
For each cut, choose the side containing a fixed point \(p\), and then delete
\(p\) from that side.  The resulting traces form an antichain in
\(T\setminus\{p\}\), because the chosen sides of two crossing cuts are
incomparable.  If \(t=2r\), Sperner's theorem therefore gives the upper bound
\(\binom{2r-1}{r-1}\), attained by the cuts whose chosen sides are the
\(r\)-subsets containing \(p\) \cite{Sperner1928}.

Now let \(t=2r+1\), and write \(U=T\setminus\{p\}\).  Crossing also implies
that \(F\cup F'\ne U\) for any two distinct traces \(F\) and \(F'\).  Hence their
complements in \(U\) form an intersecting antichain.  Milner's theorem gives
the upper bound \(\binom{2r}{r-1}\), again attained by the cuts whose chosen
sides are the \(r\)-subsets containing \(p\) \cite{Milner1968}.  Finally, the
ratios
\[
\frac{q(2r+1)}{q(2r)}=\frac{2r}{r+1}>1,
\qquad
\frac{q(2r+2)}{q(2r+1)}=\frac{2r+1}{r}>1
\]
give strict monotonicity.
\end{proof}

For a finite simple \(3\)-uniform hypergraph \(H\) and \(T\in V(H)\), let
\(L_H(T)\) denote the link graph:
\[
R\sim_{L_H(T)}S
\quad\Longleftrightarrow\quad
\{T,R,S\}\in E(H).
\]
For a bipartition \(T\), put
\[
\tau(T):=\min\{|C|:C\in T\}.
\]

\begin{lemma}[Middle-layer invariance]\label{lem:middle-invariant}
Every endomorphism of \(H_n\) maps \(B_n\) into \(B_n\).
\end{lemma}

\begin{proof}
We first prove that \(\omega(L_{H_n}(T))\le q(\tau(T))\).  The assertion is
immediate for a clique of order at most one.  For a larger clique
\(\mathcal C\), associate to each \(R\in\mathcal C\) the unordered cut
induced by \(R\) on a block of \(T\) of size \(\tau(T)\).  For
\(R,S\in\mathcal C\), the four
cells of these local cuts are precisely the four cells of the triple
\((T,R,S)\) inside that block.  Thus the local cuts cross.  The restriction
is injective, since equal unordered cuts would leave two of those cells
empty.  Therefore
\[
\omega(L_{H_n}(T))\le q(\tau(T)).
\]
If \(T\in B_n\), let \(b=\lfloor n/2\rfloor\).  On each block, choose the
extremal family from Proposition~\ref{prop:cross}, orienting every chosen
side towards its fixed point.  Retain all \(q(b)\) cuts on the
\(b\)-block and \(q(b)\) cuts on the other block; when \(n\) is odd, discard
the surplus cuts from the \(q(b+1)\)-family.  Pair the two local families
arbitrarily and take the unions of corresponding oriented sides.  Every pair
of resulting global cuts realises all four joint colours inside each block of
\(T\), and hence forms an edge with \(T\).  Each local side has size at least
\(2\), so both blocks of every global cut have size at least \(4\); the fixed
orientations also make the global cuts distinct.  We have constructed a
\(q(b)\)-clique in \(L_{H_n}(T)\), including the boundary cases \(n=8,9\).

An endomorphism maps every link clique of size at least two injectively into
the link at the image vertex.  Thus, for \(T\in B_n\),
\[
q(b)=\omega(L_{H_n}(T))
\le\omega(L_{H_n}(f(T)))
\le q(\tau(f(T))).
\]
Strict monotonicity and \(\tau(f(T))\le b\) imply \(\tau(f(T))=b\).
\end{proof}

\begin{lemma}[Separation by the balanced layer]\label{lem:separation}
For distinct \(X,Y\in V(H_n)\), there exist \(R,S\in B_n\) such that
\[
\{X,R,S\}\in E(H_n),
\qquad
\{Y,R,S\}\notin E(H_n).
\]
\end{lemma}

\begin{proof}
Write the two blocks of \(X\) as \(X_0\) and \(X_1\).  Since \(X\ne Y\), some block
\(Y_*\) of \(Y\) meets both blocks of \(X\); otherwise the two partitions
coincide.  Put \(D=Y_*^c\).

Colour \(\Omega\) with the four colours \(00,01,10,11\).  Choose one point
in each \(X_i\cap Y_*\) and colour those two points \(11\).  In each \(X_i\),
colour three additional points \(00,01,10\).  Every \(X_i\) now contains
all four colours, while \(D\) contains no point of colour \(11\).

For \(n=2k\), complete the colour multiplicities to
\[
(c_{00},c_{01},c_{10},c_{11})=(2,k-2,k-2,2);
\]
for \(n=2k+1\), complete them to
\[
(3,k-2,k-2,2).
\]
The eight points already chosen contribute two occurrences of each colour.
Since \(k\ge4\), the displayed targets dominate these initial counts and sum
to \(n\), so the remaining points can be coloured to achieve them.  Define
\[
R^*=\{x:c(x)\in\{10,11\}\},
\qquad
S^*=\{x:c(x)\in\{01,11\}\},
\qquad R=[R^*],\quad S=[S^*].
\]
The displayed multiplicities give
\(|R^*|=|S^*|=\lfloor n/2\rfloor\), so \(R\) and \(S\) belong to \(B_n\).  Every block of
\(X\) sees all four joint colours, whereas \(D\cap R^*\cap S^*=\varnothing\).
This proves the two claimed incidences.
\end{proof}

\begin{proof}[Proof of Theorem~\ref{thm:lift}]
Let \(f\in\End(H_n)\).  By Lemma~\ref{lem:middle-invariant},
\(g=f|_{B_n}\) is an endomorphism of \(B_n\), hence an automorphism.  Let
\(d\) be its finite order and set \(h=f^d\).  Then \(h\) fixes \(B_n\)
pointwise.

If \(h(X)=Y\ne X\), choose \(R,S\in B_n\) from
Lemma~\ref{lem:separation}.  Applying \(h\) would map the edge \(\{X,R,S\}\) to the
nonedge \(\{Y,R,S\}\), a contradiction.  Hence \(f^d=h=\mathrm{id}\), so
\(f\) is bijective and its inverse is the endomorphism \(f^{d-1}\).
\end{proof}

\section{Even balanced layers with
\texorpdfstring{\(n\equiv2\pmod4\)}{n congruent to 2 modulo 4}}

Let \(n=2k\), where \(k\ge5\) is odd, and let \(\widehat\Omega\) be a
\(2k\)-element ground set.  Put
\[
B_{2k}=3\text{-}\mathrm{UQI}(2k,2),
\qquad G_{2k}=[B_{2k}]_2.
\]
The vertices are unordered equipartitions \([A]=\{A,A^c\}\), with
\(|A|=k\).  If \(s=|A\cap B|\), the four intersection cells, in the order
\[
A\cap B,\quad A\setminus B,\quad B\setminus A,\quad
\widehat\Omega\setminus(A\cup B),
\]
have sizes \(s,k-s,k-s,s\).  A third equipartition, represented by a
\(k\)-set \(C\), splits all four cells if and only if
\[
2\le s\le k-2.
\]
Necessity follows because each cell must contain a point of \(C\) and a
point outside \(C\), and hence must have size at least \(2\).  For
sufficiency, choose \(C\) so that its intersections with the four cells have
respective sizes
\[
1,\quad 1,\quad k-s-1,\quad s-1.
\]

Write the ground set as
\(\widehat\Omega=\Omega\cup\{\infty\}\), with \(|\Omega|=2k-1\).  Orient
each equipartition towards the block containing \(\infty\), and delete
\(\infty\).  This identifies the vertices bijectively with the \(r\)-subsets
of \(\Omega\), where
\[
r=k-1,\qquad |\Omega|=2r+1=2k-1.
\]
If two reduced representatives meet in \(j\) points, their corresponding
oriented \(k\)-blocks meet in \(j+1\) points.  Thus distinct vertices are
nonadjacent in \(G_{2k}\) precisely when \(j\in\{0,r-1\}\).  Let \(\Gamma\)
be the graph on \(\binom{\Omega}{r}\) in which two sets are adjacent exactly
in those two cases.
Then
\[
G_{2k}=\overline{\Gamma}.
\]

\begin{lemma}[Clique number of \(\Gamma\)]\label{lem:gamma-clique}
\(\Gamma\) is vertex-transitive and \(\omega(\Gamma)=r+2\).
\end{lemma}

\begin{proof}
The natural action of \(S_{\Omega}\) on \(r\)-subsets is transitive and
preserves \(\Gamma\).  For the lower bound, fixing an \((r-1)\)-subset gives
a clique of size \(r+2\).  For the upper
bound, first suppose that a clique contains disjoint \(A,B\).  There is one
point \(x\) outside \(A\cup B\).  If \(C\) is adjacent to both, then
\[
|C\cap A|+|C\cap B|+\boldsymbol{1}_{x\in C}=r,
\]
with both intersections in \(\{0,r-1\}\).  Hence \(C\) is
\[
(A\setminus\{a\})\cup\{x\}
\quad\text{or}\quad
(B\setminus\{b\})\cup\{x\}.
\]
One set of each type meets in the singleton \(\{x\}\), so the two types
cannot be mixed when \(r\ge3\).  The clique has at most \(2+r=r+2\)
members.

Now suppose that there is no disjoint pair.  All pairwise intersections
have size \(r-1\).  Write
\[
A=D\cup\{a\},\qquad B=D\cup\{b\},\qquad |D|=r-1.
\]
Indeed, a set \(C\) meeting both \(A\) and \(B\) in \(r-1\) points either
contains \(D\), or omits some \(d\in D\).  In the latter case the two
intersection conditions force
\(C=(D\setminus\{d\})\cup\{a,b\}\subseteq D\cup\{a,b\}\).
Every further member either contains \(D\), or is an \(r\)-subset of
\(D\cup\{a,b\}\).  If only the first type occurs, the family is contained
in the \(r+2\) sets obtained by adjoining one point to \(D\).  If a member of
the second type occurs, it is not adjacent in \(\Gamma\) to any set
\(D\cup\{z\}\) with \(z\notin D\cup\{a,b\}\), and the clique is contained in
the \(r+1\) \(r\)-subsets of \(D\cup\{a,b\}\).  This proves the bound.
\end{proof}

Recall that a Steiner system \(S(t,k,v)\) is a family of \(k\)-subsets of a
\(v\)-set in which every \(t\)-subset lies in exactly one block.

\begin{lemma}[Design implication]\label{lem:design-implication}
If
\[
\chi(\overline{\Gamma})=\omega(\overline{\Gamma}),
\]
then a Steiner system \(S(r-1,r,2r+1)\) exists.
\end{lemma}

\begin{proof}
Let \(\mathcal Q\) be a clique and \(\mathcal I\) an independent set of the
vertex-transitive graph \(\Gamma\).  Averaging
\(|g\mathcal Q\cap\mathcal I|\le1\) over a transitive automorphism group gives
\[
|\mathcal Q||\mathcal I|\le\binom{2r+1}{r}.
\]
Assume \(\chi(\overline{\Gamma})=\omega(\overline{\Gamma})=c\), and put
\(N=\binom{2r+1}{r}\).  Since a colour class of \(\overline{\Gamma}\) is a
clique of \(\Gamma\), Lemma~\ref{lem:gamma-clique} gives
\(N\le c(r+2)\).  On the other hand,
\(c=\omega(\overline{\Gamma})=\alpha(\Gamma)\), so the clique--coclique
inequality above gives \(c(r+2)\le N\).  Therefore equality holds
throughout:
\[
c=\frac{\binom{2r+1}{r}}{r+2}
\]
and every colour class has size \(r+2\).

Take a maximum clique \(\mathcal B\) of \(\overline{\Gamma}\).  It is an
independent set in \(\Gamma\), so no two of its \(r\)-subsets share an
\((r-1)\)-subset.  Moreover,
\[
|\mathcal B|\,r
=\frac{\binom{2r+1}{r}}{r+2}\,r
=\binom{2r+1}{r-1}.
\]
Thus every \((r-1)\)-subset lies in exactly one member of \(\mathcal B\), so
\(\mathcal B\) is the block set of a Steiner system \(S(r-1,r,2r+1)\).
\end{proof}

\begin{proposition}\label{prop:mod2}
If \(k\ge5\) is odd, then \(3\text{-}\mathrm{UQI}(2k,2)\) is a core.
\end{proposition}

\begin{proof}
The group \(S_{\widehat\Omega}\cong S_{2k}\) acts on unordered
equipartitions and preserves \(G_{2k}\).  By the adjacency criterion above,
a distinct nonedge has
\(s\in\{1,k-1\}\); the cases \(s=0\) and \(s=k\) represent the same unordered
equipartition.  After choosing orientations, the four intersection cells
therefore have sizes
\(1,k-1,k-1,1\).  The alternatives \(s=1\) and \(s=k-1\) describe the same
unordered pair after complementing one chosen block.  Hence
\(S_{\widehat\Omega}\), and
therefore \(\Aut(G_{2k})\), is transitive on unordered nonedges.
Corollary 2.2 of Cameron and Kazanidis implies that \(G_{2k}\) is either a
core or a graph with a complete core \cite{CameronKazanidis2008}.  A finite
graph has a complete core exactly when its clique and chromatic numbers agree.

By Lemma~\ref{lem:design-implication}, equality of the clique and chromatic
numbers would yield
\[
S(k-2,k-1,2k-1).
\]
In such a design, the number of blocks through a fixed \((k-3)\)-subset is
\[
\lambda_{k-3}
=\frac{\binom{k+2}{1}}{\binom21}
=\frac{k+2}{2},
\]
which is not an integer for odd \(k\).  Hence the complete-core alternative
is impossible, so \(G_{2k}\) is a core.  Lemma~\ref{lem:transfer} now shows
that \(B_{2k}\) is a core.
\end{proof}

\section{Even balanced layers with
\texorpdfstring{\(n\equiv0\pmod4\)}{n congruent to 0 modulo 4}}

Let \(n=4m\) and put \(B_{4m}=3\text{-}\mathrm{UQI}(4m,2)\).  Write the
\(4m\)-element ground set as
\(\widehat\Omega=\Omega\cup\{\infty\}\), where \(|\Omega|=4m-1\).  Orient
each equipartition towards its unique block containing \(\infty\).
Removing \(\infty\) from that block gives a bijection with the fixed-point
model
\[
V(B_{4m})=\binom{\Omega}{2m-1},
\qquad |\Omega|=4m-1.
\]
Here, the bit \(1\) denotes membership in the corresponding reduced set,
equivalently membership in its oriented block containing \(\infty\).
For three vertices, the \(111\)-cell contains \(\infty\) automatically, so
they form an edge precisely when the other seven Boolean cells in \(\Omega\)
are nonempty.

Put \(r=2m-1\) and \(a=|X\cap Y|\).  In the order
\[
X\cap Y,\quad X\setminus Y,\quad Y\setminus X,\quad
\Omega\setminus(X\cup Y),
\]
the four pair regions have sizes \(a,r-a,r-a,a+1\).  If
\(1\le a\le r-2\), choose a third \(r\)-set whose intersections with these
four regions have respective sizes
\[
0,\quad 1,\quad r-a-1,\quad a.
\]
The first region need only contribute to its
non-\(111\) cell, whereas each of the other three regions is split between
the two values of the third bit.  This makes all seven required cells
nonempty.  Conversely, the non-\(111\) cell in the first region and the two
cells in each difference region force \(a\ge1\) and \(r-a\ge2\).  Thus the
\(2\)-section joins \(X\) and \(Y\) exactly when
\[
1\le |X\cap Y|\le2m-3.
\]
Distinct nonedges therefore have intersection \(0\) or \(r-1\).  In the
equipartition model these belong to the same orbit: complementing one
oriented block interchanges the two cases, and the four cells have sizes
\(1,r,r,1\).  Hence all unordered nonedges form one orbit under the action of
\(S_{\widehat\Omega}\cong S_{4m}\) transported to the fixed-point model.

\begin{lemma}[A collision criterion]\label{lem:critical-fold}
Suppose that \(H\) is a vertex-transitive \(3\)-uniform hypergraph, every
link of \(H\) has clique number at most \(M\), and \(\Aut(H)\) is transitive
on the unordered pairs of distinct vertices that are nonedges of \([H]_2\).
Suppose
also that a fixed nonedge \(\{P,Q\}\) admits distinct vertices
\(R_1,\ldots,R_{M+1}\) such that, for every \(i<j\), at least one of
\[
\{P,R_i,R_j\},\qquad \{Q,R_i,R_j\}
\]
is an edge.  Then \(H\) is a core.
\end{lemma}

\begin{proof}
If an endomorphism identifies \(P\) and \(Q\), the images of
\(R_1,\ldots,R_{M+1}\) are distinct: for each \(i<j\), the pair
\(\{R_i,R_j\}\) is contained in a hyperedge and hence its vertices cannot
have the same image.  They form an
\((M+1)\)-clique in the link of the common image, a contradiction.  More
generally, two distinct vertices identified by an endomorphism must be
nonadjacent in the \(2\)-section.  By the orbit hypothesis an automorphism
sends the unordered pair \(\{P,Q\}\) to that pair, and precomposition reduces
the collision to the fixed pair \(\{P,Q\}\).  Thus no collision is possible.
\end{proof}

\begin{proposition}\label{prop:4m}
For every \(m\ge4\), \(3\text{-}\mathrm{UQI}(4m,2)\) is a core.
\end{proposition}

\begin{proof}
The action of \(S_{\widehat\Omega}\) is transitive on \(V(B_{4m})\).  For a
clique \(\mathcal C\) in \(L_{B_{4m}}(T)\), the traces
\(\{T\cap R:R\in\mathcal C\}\) form an antichain in \(2^T\): a containment
between two traces would make one of the two difference cells inside \(T\)
empty.  Hence every link of \(B_{4m}\) has clique number at most
\[
M_m:=\binom{2m-1}{m-1}.
\]

\emph{Construction.}
We construct \(M_m\) central vertices indexed by the \((m-1)\)-subsets of
\(Q\), and then add one exceptional vertex.  The three types below are
chosen so that pairs of central vertices are witnessed by \(Q\), whereas
pairs containing the exceptional vertex are witnessed by either \(P\) or
\(Q\).

Choose pairwise disjoint sets \(K\), \(\{p,q\}\), and \(W\) with
\[
|K|=2m-2,\qquad |W|=2m-1,
\]
and put \(P=K\cup\{p\}\), \(Q=K\cup\{q\}\).  Fix
\(x\in K\) and distinct \(u,v\in W\), and put \(O=\{p\}\cup W\).
The pair \(\{P,Q\}\) is a nonedge because \(|P\cap Q|=2m-2\).

Divide the family \(\binom{Q}{m-1}\) into three types, labelled
\(0,10,11\), respectively:
\[
\begin{array}{c|l}
0&x\notin A,\\
10&x\in A,\ q\notin A,\\
11&x,q\in A.
\end{array}
\]
The complement pairs of \(m\)-subsets of \(O\) likewise fall into three
classes: either \(u\) and \(v\) are separated, or they lie together in the
unique \(p\)-free representative, or they lie together in its complement.  Thus
these classes exhaust all complement pairs.  Their cardinalities are
\[
\begin{array}{c|c|c}
\text{type}&\text{condition}&\text{number}\\ \hline
0&u,v\text{ separated}&\binom{2m-2}{m-1}\\
10&\text{\(p\)-free representative contains }u,v&
\binom{2m-3}{m-2}\\
11&\text{\(p\)-free representative avoids }u,v&
\binom{2m-3}{m-3}.
\end{array}
\]
The three subfamilies of \(\binom{Q}{m-1}\) have, respectively, the three
cardinalities displayed above: one chooses \(m-1\) points while avoiding
\(x\), chooses \(x\) and \(m-2\) further points while avoiding \(q\), or
chooses \(x,q\) and \(m-3\) further points.  For each type, fix a bijection
between the two corresponding classes.

For types \(10\) and \(11\), let \(C_A\) be the unique \(p\)-free
representative of the assigned complement pair; for type \(0\), choose
either representative.  Set \(R_A=A\cup C_A\).  Since \(Q\cap O=\varnothing\),
each \(R_A\) has size \((m-1)+m=2m-1\) and hence is a vertex of \(B_{4m}\).

\emph{Pairs of central vertices.}
For distinct \(A\) and \(A'\), put
\[
s=|A\cap A'|,\qquad t=|C_A\cap C_{A'}|.
\]
Since \(A\) and \(A'\) are distinct and have the same size,
\(0\le s\le m-2\).  Since their assigned complement pairs are distinct,
\(C_A\) and \(C_{A'}\) are neither equal nor complementary, and
\(1\le t\le m-1\).  With the non-\(111\) Boolean cells indexed in the order
\(000,001,010,011,100,101,110\), their sizes for
\((Q,R_A,R_{A'})\) are
\[
\begin{array}{c|ccccccc}
 &000&001&010&011&100&101&110\\ \hline
\text{size}
 &t&m-t&m-t&t&s+1&m-1-s&m-1-s.
\end{array}
\]
They are all positive.  Hence, for distinct \(A,A'\),
\[
\{Q,R_A,R_{A'}\}\in E(B_{4m}).
\]

\emph{Pairs containing the exceptional vertex.}
Put
\[
Z=\{x,p\}\cup(W\setminus\{u,v\}).
\]
For \(A\) of a given type, let
\(\operatorname{anc}(A)\in\{P,Q\}\) be the anchor shown in the corresponding
row below.  The seven non-\(111\) Boolean-cell sizes for
\((\operatorname{anc}(A),Z,R_A)\) are
\[
\begin{array}{c|c|ccccccc}
\text{type}&\text{anchor}&000&001&010&011&100&101&110\\ \hline
0&Q&1&1&m-1&m-1&m-1&m-1&1\\
10&P&1&2&m-1&m-2&m-1&m-2&1\\
11&P&2&1&m-3&m&m&m-3&1.
\end{array}
\]
All are positive for \(m\ge4\), and hence
\(\{\operatorname{anc}(A),Z,R_A\}\in E(B_{4m})\).

The family
\[
\{Z\}\cup\left\{R_A:A\in\binom{Q}{m-1}\right\}
\]
has \(M_m+1\) vertices.  Every pair of central vertices is witnessed by
\(Q\), and every pair consisting of \(Z\) and a central vertex is witnessed
by the anchor in the table.  The hypotheses of Lemma~\ref{lem:critical-fold}
are therefore satisfied, and \(B_{4m}\) is a core.
\end{proof}

\subsection{The exceptional value
\texorpdfstring{\(n=12\)}{n=12}}\label{subsec:n12}

When \(m=3\), the two entries \(m-3\) in the type \(11\) row vanish.  We
instead use an obstruction with eleven auxiliary vertices.  After deleting
\(\infty\), let
\[
\Omega_{12}^{\circ}:=\{0,1,\ldots,10\},
\]
and use the model \(V(B_{12})=\binom{\Omega_{12}^{\circ}}5\).  Set
\[
P=\{0,1,2,3,4\},
\qquad
Q=\{0,1,2,3,5\}.
\]
The eleven auxiliary vertices and their anchor table are given in
Appendix~\ref{app:n12}; the accompanying program
\texttt{supplementary/verify\_certificates.py} verifies all \(55\) required
pairwise incidences.  Every link \(L_{B_{12}}(T)\) has clique number at most
\(\binom{5}{2}=10\), while identifying \(P\) and \(Q\) would force an
\(11\)-clique in one link.  The preceding vertex-transitivity argument and
the single-orbit property for nonedges also apply when \(m=3\).  Hence
Lemma~\ref{lem:critical-fold} applies to this certificate and proves that
\(B_{12}=3\text{-}\mathrm{UQI}(12,2)\) is a core.

\begin{proof}[Proof of Theorem~\ref{thm:even-middle}]
For \(n=8\), all admissible vertices have block sizes \(4+4\), so
\(3\text{-}\mathrm{QI}(8,2)=3\text{-}\mathrm{UQI}(8,2)\); this case is due
to Thomas and Akhtar \cite[Theorem 4]{ThomasAkhtar2026}.
Proposition~\ref{prop:mod2} treats
\(n\equiv2\pmod4\), the preceding subsection treats \(n=12\), and
Proposition~\ref{prop:4m} treats \(n\equiv0\pmod4\), \(n\ge16\).
\end{proof}

\begin{proof}[Proof of Theorem~\ref{thm:main}]
For even \(n\), the balanced layer is a core by
Theorem~\ref{thm:even-middle}.  For odd \(n\ge9\), the balanced layer is a
core by Theorem 7.1 and Corollary 7.2 of the companion manuscript
\cite{KimMergedJohnson2026}.  Theorem~\ref{thm:lift} therefore proves the
result.  The case \(n=8\) is already included in the even statement.
\end{proof}

\appendix

\section{The \texorpdfstring{\(n=12\)}{n=12} collision certificate}\label{app:n12}

The auxiliary vertices \(R_1,\ldots,R_{11}\) are
\[
\begin{array}{c|l@{\qquad}c|l}
R_1&\{1,2,3,5,10\}&R_2&\{0,3,8,9,10\}\\
R_3&\{0,2,7,9,10\}&R_4&\{2,3,4,9,10\}\\
R_5&\{1,3,4,8,10\}&R_6&\{1,2,4,7,10\}\\
R_7&\{0,5,7,8,9\}&R_8&\{0,1,6,7,10\}\\
R_9&\{3,4,5,7,9\}&R_{10}&\{2,4,5,6,9\}\\
R_{11}&\{1,4,5,7,8\}&&
\end{array}
\]
and the upper-triangular anchor table is
\[
\begin{array}{c|cccccccccc}
 &2&3&4&5&6&7&8&9&10&11\\ \hline
1&P&P&P&P&P&P&P&P&P&P\\
2& &P&Q&Q&Q&Q&P&P&P&P\\
3& & &Q&Q&Q&Q&P&P&P&P\\
4& & & &P&P&P&Q&Q&Q&Q\\
5& & & & &P&P&P&Q&Q&Q\\
6& & & & & &P&Q&P&Q&Q\\
7& & & & & & &Q&Q&P&Q\\
8& & & & & & & &P&P&P\\
9& & & & & & & & &P&P\\
10& & & & & & & & & &P
\end{array}
\]
An entry \(P\) or \(Q\) indicates that the corresponding pair of auxiliary
vertices forms a hyperedge with that anchor.  Each incidence is a direct
seven-cell check; for example, the entry in row \(1\), column \(2\) states
that, in the order \(000,001,010,011,100,101,110\), the seven required cell
sizes for \((P,R_1,R_2)\) are \((2,2,1,1,1,1,2)\).  The accompanying script
verifies all \(55\) incidences directly from the listed sets.

\section*{Declaration on generative AI}

ChatGPT assisted with exploratory calculations, code development, drafting,
and editorial review.  The author is responsible for the mathematical
statements, proofs, and accompanying code.

\section*{Data and code availability}

The source archive accompanying this preprint contains the verification
scripts.  They check the finite \(n=12\) certificate, the boundary instance
\(n=16\) of the general construction, and the stated small lift instances.

\bibliographystyle{plain}
\bibliography{references}

\end{document}
